\section{The Clock Genesis Prints}

Genesis 15 opens ``after these things'' with a divine word to Abram ``in a
vision.'' Abram objects that he remains childless. God answers that an heir
from his own body will inherit, brings him outside, commands him to look and
count, and promises offspring like the stars. Abram believes. The chapter then
moves to the land promise, the animals, the birds, sunset, deep sleep, darkness,
oracle, fire, and covenant.

Three temporal statements control the immediate issue:

\begin{itemize}
\item At verse 12, the sun is going down or about to set.
\item At verse 17, the sun has set and darkness has come.
\item At verse 18, the covenant is made ``in that day.''
\end{itemize}

No explicit night precedes verse 12. No formula says ``on the next day'' between
the star promise and the animals. Read in narrated order as one continuous
scene, verse 5 is before the chapter's sunset. That is the daylight argument in
its strongest form.

It is not a grammatical proof. Biblical narrative can compress time and omit
transitions. ``After these things'' can open a new episode without measuring
the interval from the preceding one. A vision can contain represented time
that does not map straightforwardly onto waking chronology. Continuity is the
natural surface reading, not the only syntactically possible reconstruction.

